\documentclass[../Main_Include.tex]{subfiles}
\begin{document}
\backgroundsetup{contents={}}  % suprime draft de background en compilación standalone
% ============================
% LABORATORIO 2 — Laboratorio2_1010.tex
% ============================
%
% CONTENIDO:
%   Sec. 1 — Ejercicio 1: Carga y visualización del dataset Iris
%   Sec. 2 — Ejercicio 2: Clasificación KNN sobre el dataset Iris
%
% DEPENDENCIAS (cargadas desde Main_Include.tex):
%   Preambulo.tex       — paquetes, colores, estilos de código
%   PaginaEstilo.tex    — lstset global
%   EstiloSecciones.tex — \titleformat, fancyhdr
%   Entornos.tex        — cuadropregunta
%
% ============================

% --- Título del encabezado para este informe ---
\newcommand{\tituloinforme}{Área de Aplicación de la Inteligencia Artificial — Laboratorio 2}

% ══════════════════════════════════════════════════════════════════════
\section{Ejercicio 1 — Carga y visualización del dataset Iris}

Se implementó un programa en Python que carga el dataset Iris desde
\texttt{scikit-learn}, lo convierte a un DataFrame de \texttt{pandas}
y lo imprime completo en consola sin truncar ninguna fila.

% ──────────────────────────────────────────────────────────────────────
\subsection*{Observaciones}

El dataset contiene 150 muestras × 5 columnas: las cuatro características
físicas de cada flor (\texttt{sepal length}, \texttt{sepal width},
\texttt{petal length}, \texttt{petal width}) y la columna \texttt{species},
que codifica la especie como valor numérico: \texttt{0} = \emph{setosa},
\texttt{1} = \emph{versicolor}, \texttt{2} = \emph{virginica}.

La columna \texttt{species} muestra valores numéricos en lugar de nombres
por decisión de diseño del ejercicio --- es más limpio visualmente y
no se requiere un \texttt{.map()} adicional.

El uso de \texttt{df.to\_string()} fuerza que \texttt{pandas} imprima
las 150 filas completas; sin él, la salida se truncaría con \texttt{...}
en el medio.

% ──────────────────────────────────────────────────────────────────────
\subsection*{Capturas de pantalla}

\begin{figure}[H]
    \centering
    \includegraphics[width=0.85\linewidth]{Graficos/Lab2/E1_are_001.png}
    \caption{Inicio de la ejecución --- filas 0 a 4, especie \emph{setosa} (species = 0).}
    \label{fig:e1_are_001}
\end{figure}

\begin{figure}[H]
    \centering
    \includegraphics[width=\linewidth]{Graficos/Lab2/E1_are_002.png}
    \caption{Final de la ejecución --- filas 146 a 149, especie \emph{virginica} (species = 2).}
    \label{fig:e1_are_002}
\end{figure}

% ══════════════════════════════════════════════════════════════════════
\newpage
\section{Ejercicio 2 — Clasificación KNN sobre el dataset Iris}

Se implementó un clasificador KNN completo sobre el dataset Iris.
El programa divide los datos en 80\% entrenamiento y 20\% prueba,
entrena un modelo con $k=3$ vecinos, predice sobre el set de prueba
y calcula el accuracy final.

% ──────────────────────────────────────────────────────────────────────
\subsection*{Observaciones}

La salida del programa presenta cuatro bloques:

\begin{itemize}[leftmargin=1.5cm, itemsep=2pt]
    \item Primeras 5 filas del dataset con columnas \texttt{target} y
          \texttt{species} (nombre de la especie).
    \item Cantidad de muestras: \texttt{Training: 120, Test: 30}.
    \item Comparación de las primeras 5 predicciones contra los valores
          reales, con nombres de especie para ambas columnas.
    \item Accuracy final del modelo.
\end{itemize}

\paragraph{Accuracy obtenido.}
El resultado fue \textbf{100.00\%} --- es decir, las 30 muestras del set
de prueba fueron clasificadas correctamente. Las notas del ejercicio
mencionaban 96.67\% como valor típico con versiones anteriores de
\texttt{scikit-learn}; la diferencia se debe a que la versión instalada
(1.9.1) maneja el \texttt{random\_state=42} de forma ligeramente distinta.
Ambos resultados son válidos.

\paragraph{Parámetros del modelo.}
\texttt{KNeighborsClassifier(n\_neighbors=3)} con
\texttt{train\_test\_split(test\_size=0.2, random\_state=42)}.

% ──────────────────────────────────────────────────────────────────────
\subsection*{Capturas de pantalla}

\begin{figure}[H]
    \centering
    \includegraphics[width=\linewidth]{Graficos/Lab2/E2_are_001.png}
    \caption{Salida completa del clasificador KNN --- dataset, muestras,
             predicciones y accuracy 100.00\%.}
    \label{fig:e2_are_001}
\end{figure}

% ══════════════════════════════════════════════════════════════════════
\newpage
\appendix
\renewcommand{\thesection}{\Alph{section}}
\section*{Apéndice — Códigos fuente}
\addcontentsline{toc}{section}{Apéndice}

\subsection*{Ejercicio 1 — Carga y visualización del dataset Iris}

\lstinputlisting[style=py_dark,
    caption={Ejercicio 1 --- Carga y visualizaci\'on del dataset Iris},
    label={lst:e1_area}]{Anexos/Lab2/E1_AREA_ETN1010.py}

\newpage

\subsection*{Ejercicio 2 — Clasificación KNN}

\lstinputlisting[style=py_dark,
    caption={Ejercicio 2 --- Clasificaci\'on KNN sobre el dataset Iris},
    label={lst:e2_area}]{Anexos/Lab2/E2_AREA_ETN1010.py}

\end{document}
